\chapter{Purpose and Scope}

\section{Purpose of the Project}

Enterprise connector delivery in telecom teams was historically fragmented across manual specification reading, JAR inspection, mapping spreadsheets, YAML editing, and script execution. This process produced inconsistent outputs and weak traceability.

The API Development Tool addresses that gap by consolidating the full lifecycle in one guided web application. A React 18 frontend~\cite{react_docs} drives the workflow, and a Node.js/Express backend~\cite{express_docs} manages parsing, persistence, synchronisation, and generation orchestration. The purpose is operationally clear: shorten delivery time, standardise mapping outputs, and preserve an auditable record of all connector decisions.

\section{Functional Scope}

The system scope covers five sequential phases that transform input artifacts into a deployable connector module. Figure~\ref{fig:workflow_lifecycle} shows the artifact flow.

\begin{figure}[H]
\centering
\begin{tikzpicture}[
    phase/.style={rectangle, rounded corners=6pt, line width=0.7pt,
                  text width=10cm, align=center, minimum height=1.5cm,
                  inner xsep=10pt, inner ysep=8pt,
                  drop shadow={shadow xshift=0.8pt, shadow yshift=-0.8pt, opacity=0.2}},
    badge/.style={circle, minimum size=0.7cm, font=\footnotesize\bfseries,
                  text=white, inner sep=0pt, line width=0.7pt,
                  drop shadow={shadow xshift=0.5pt, shadow yshift=-0.5pt, opacity=0.15}},
    arr/.style={->, line width=1.3pt, >=stealth, color=black!40},
    node distance=1.2cm
]

% Phase 1 - Blue
\node[phase, draw=blue!55!black, fill=blue!12] (p1)
    {\textbf{\small Module Management}\\[3pt]
     \footnotesize Create, select, or delete connector modules; normalise names; establish directory structure};
\node[badge, fill=blue!60!black, left=0.45cm of p1.west] {\textsf{1}};

% Phase 2 - Teal
\node[phase, draw=teal!55!black, fill=teal!10, below=of p1] (p2)
    {\textbf{\small Specification Ingestion \& Parsing}\\[3pt]
     \footnotesize Upload Swagger/OpenAPI YAML; resolve \$ref pointers; flatten schemas; extract operation metadata};
\node[badge, fill=teal!60!black, left=0.45cm of p2.west] {\textsf{2}};

% Phase 3 - Green
\node[phase, draw=green!45!black, fill=green!8, below=of p2] (p3)
    {\textbf{\small JAR Introspection (EJB Metadata Extraction)}\\[3pt]
     \footnotesize Read bytecode from client-kit JARs; detect EJB interfaces; extract method signatures and field paths};
\node[badge, fill=green!50!black, left=0.45cm of p3.west] {\textsf{3}};

% Phase 4 - Amber
\node[phase, draw=orange!55!black, fill=orange!10, below=of p3] (p4)
    {\textbf{\small Interactive Mapping \& Configuration Generation}\\[3pt]
     \footnotesize Map API parameters to EJB fields via dropdowns; apply source inference; serialise to YAML config};
\node[badge, fill=orange!60!black, left=0.45cm of p4.west] {\textsf{4}};

% Phase 5 - Red
\node[phase, draw=red!50!black, fill=red!8, below=of p4] (p5)
    {\textbf{\small Module Generation \& Build Execution}\\[3pt]
     \footnotesize Invoke generation scripts; scaffold Maven module; stream real-time output via Socket.IO};
\node[badge, fill=red!55!black, left=0.45cm of p5.west] {\textsf{5}};

% Arrows with colour-matched labels
\draw[arr] (p1.south) -- node[right=5pt, font=\scriptsize, fill=blue!6,
    draw=blue!30, rounded corners=3pt, inner sep=4pt, text=black!80]
    {Module context + directory path} (p2.north);
\draw[arr] (p2.south) -- node[right=5pt, font=\scriptsize, fill=teal!6,
    draw=teal!30, rounded corners=3pt, inner sep=4pt, text=black!80]
    {Parsed API operations + field paths} (p3.north);
\draw[arr] (p3.south) -- node[right=5pt, font=\scriptsize, fill=green!6,
    draw=green!25, rounded corners=3pt, inner sep=4pt, text=black!80]
    {EJB method signatures + flattened field paths} (p4.north);
\draw[arr] (p4.south) -- node[right=5pt, font=\scriptsize, fill=orange!6,
    draw=orange!30, rounded corners=3pt, inner sep=4pt, text=black!80]
    {User-config YAML (endpoint defs, call chains, mappings)} (p5.north);
\end{tikzpicture}
\caption{Five-phase workflow lifecycle of the API Development Tool, showing the key artifacts passed between each phase.}
\label{fig:workflow_lifecycle}
\end{figure}

The phases are:
\begin{enumerate}
    \item \textbf{Module management}: create/select/delete modules and enforce normalised names plus directory conventions.
    \item \textbf{Specification ingestion}: upload Swagger/OpenAPI YAML, resolve references, flatten schemas, and extract operation metadata.
    \item \textbf{JAR introspection}: parse compiled client-kit bytecode, detect EJB interfaces, and extract method and field-path metadata with SHA-256 deduplication.
    \item \textbf{Interactive mapping \& config generation}: map request/response fields (including multi-EJB chains), apply source inference, and generate YAML.
    \item \textbf{Module generation}: invoke the PowerShell-driven generation flow and stream execution output through Socket.IO.
\end{enumerate}

\section{Persistence and Synchronisation}

State is maintained through a hybrid model: sql.js in-process SQLite~\cite{sqljs_docs} for metadata and filesystem storage for artifacts. The core tables are \texttt{swagger\_files}, \texttt{client\_kit\_metadata}, and \texttt{user\_configs}, with unique constraints for duplicate prevention.

Synchronization combines startup reconciliation, runtime chokidar watchers~\cite{chokidar_docs}, and frontend polling fallback. Reconciliation repairs drift between disk and database; watchers detect external file changes; Socket.IO broadcasts keep clients aligned.

\section{Non-Functional Considerations}

The platform supports HTTPS/HTTP modes. In HTTPS mode, it auto-generates self-signed certificates at first startup, which fits air-gapped deployment constraints. The frontend is packaged as a Tomcat WAR~\cite{tomcat_docs}; the backend runs under PM2.

Portability is achieved through relative path resolution and dependency bundling. The stack runs without external database services or runtime internet access. A PowerShell packaging script automates validation, frontend build, WAR creation, and deployable ZIP assembly.

\section{Future Enhancement Opportunities}

The current scope covers the full connector workflow. Planned extensions include authentication and role control, cross-platform generation, stronger automated testing, containerised deployment, and multi-tenant isolation. The existing modular design can support these additions without architectural reset.

